\printedpage{14}
% head: A. Grothendieck : Dérivateurs, Ch. VII, Édition des Universités de Montpellier II et Paris VII
% foot: Version 22 décembre 2001

by simply taking the ``law'' corresponding to $\varphi$
\[
c \longmapsto c\circ\varphi \qquad \text{for } c : I_\tau \longrightarrow X.
\]
More precisely, the natural map
\[
\underbrace{\Hom_{\Cat}(I_{\tau'},I_\tau)}_{=\mathrm{Ob}\,\uHom(I_{\tau'},I_\tau)} \longrightarrow \Hom(\mathrm{Ch}_\tau,\mathrm{Ch}_{\tau'}),
\]
\msp{page 22}
where the second Hom is regarded as a Hom of objects in $\Cat^{\vee}=\uHom(\Cat,\mathrm{Ens})$. Thus, in the two examples above, where $\tau=\bDelta_2^+$ and $\tau'=\bDelta_1^+$ or $\bDelta_2^+$, the functors in the reverse direction are
\[
\begin{tikzcd}
\bDelta_2^+\ \Big\{ & 0 \arrow[r] & 1 \arrow[r] & 2 \\
\bDelta_1^+\ \Big\{ & 0 \arrow[u, mapsto] \arrow[rr] & & 1 \arrow[u, mapsto]
\end{tikzcd}
\]
and
\[
\begin{tikzcd}
\bDelta_2^+\ \Big\{ & 0 \arrow[r] & 1 \arrow[r] & 2 \\
\bDelta_2^+\ \Big\{ & 0 \arrow[u, mapsto] \arrow[r] & 1 \arrow[ur, mapsto] \arrow[r] & 2\,. \arrow[u, mapsto]
\end{tikzcd}
\]
This is a general principle for reducing all the $\pm$ canonical ``operations'' on paths in arbitrary categories, and statements concerning them, to questions where $X$ disappears and where we work in the interval categories themselves. One of the technical advantages is that these are very small categories, giving rise moreover to a sure topological intuition (which risks evaporating at times when working in categories of paths), and also that in such categories, hence also in the categories of functors between them, there is always at most
\msp{page 23}
\textit{one single arrow} going from a given object to another. In other words, in such categories all diagrams are commutative.

Of course, we also see that if we have a law of formation $L$ of paths of type $\tau'$ from paths of type $\tau$ (with the condition of functoriality with respect to the category $X$ in which we take the paths), this law is necessarily \textit{functorial with respect to paths:} if we have a homomorphism
\[
c \xrightarrow{v} c' \qquad \text{in } \uCh_\tau(X),
\]
\printedpage{15}
% head: A. Grothendieck : Dérivateurs, Ch. VII, Édition des Universités de Montpellier II et Paris VII
% foot: Version 22 décembre 2001
then the law $L$ ``extends'' to such $v$ to give
\[
L(c) \xrightarrow{L(v)} L(c'),
\]
with the usual relations
\[
L(\mathrm{id}_C)=\mathrm{id}_{L(c)}, \qquad L(v'v)=L(v')L(v). % [sic] \mathrm{id}_C -> \mathrm{id}_c, capital C printed while the chemin is c (cf. L(c))
\]
This follows from the description of $L$ as an operation $L_\varphi$ in terms of a well-determined functor
\[
I_{\tau'} \xrightarrow{\varphi} I_\tau,
\]
which in turn gives for every $X$ a \textit{functor}
\[
\left\{
\begin{array}{rcl}
\uCh_\tau(X) & \longrightarrow & \uCh_{\tau'}(X)\\[1ex]
c & \longmapsto & c\circ\varphi
\end{array}
\right.
\]
more precise than simply a \textit{map} $\mathrm{Ch}_\tau(X)\longrightarrow\mathrm{Ch}_{\tau'}(X)$ between the corresponding sets of objects. % [sic] un application -> une application, gender agreement

\msp{page 24}
We now have the choice between

\begin{enumerate}[label=\alph*)]
\item A more or less systematic study of functors between intervals. I prefer to wait for this until I need this and that, and to do only the strict minimum for the constructions I have in view.
\item A description of the compositions of paths. % [sic] la compositions -> la composition, number agreement
\item A formalism of categories of paths \textit{of unspecified type.}
\end{enumerate}

Since my aim is above all to describe such categories $\uCh(X)$ (without a subscript $\tau$!), with good properties, I will start with c), then look at b) as a crucial test for the usability of the notions obtained, and finally take up a) when I know what I will need of it.

\printedpage{16}
% head: A. Grothendieck : Dérivateurs, Ch. VII, Édition des Universités de Montpellier II et Paris VII
% foot: Version 22 décembre 2001

\section*{4\quad Category of paths $\uCh(X)$}

I want to ``include'' all the categories $\uCh_\tau(X)$ of paths of all the types one wants, in one and the same category $\uCh(X)$. Recall
\[
\uCh_\tau(X)=\uHom(I_\tau,X),
\]
where $I_\tau$ is the standard (oriented) interval of type $\tau$.

The set of all types of intervals
\msp{page 25}
can be identified with the free monoid Ty on two generators $\Delta^+$ and $\Delta^-$, whose elements are arbitrary finite sequences (possibly empty) of objects alternately equal to $\Delta^+$ and to $\Delta^-$. If the first element is $\Delta^+$, we have a type of path with direct initial direction, retrograde otherwise. The decomposition of $I_\tau$, or of $\mathrm{Fl}\,(I_\tau)$, into critical intervals corresponds to the reduced expression of a word as % [sic] eventuellement -> éventuellement, missing accent
\[
(\Delta^+)^{l_1}(\Delta^-)^{l_2}\ \cdots\ (\Delta^?)^{l_r} \qquad (\Delta^?=\Delta^+ \text{ if } r \text{ odd, } \Delta^?=\Delta^- \text{ if } r \text{ even}),
\]
or as
\[
(\Delta^-)^{l_1}(\Delta^+)^{l_2}\ \cdots\ (\Delta^?)^{l_r} \qquad (\Delta^?=\Delta^- \text{ if } r \text{ odd, } \Delta^?=\Delta^+ \text{ if } r \text{ even}).
\]
The $l_i$ are the lengths of the successive critical intervals (or of the successive structural chains), their number is the reduced length of the interval.

We are going to regard Ty as the set of objects of a category $\uTy$:
\[
\left\{
\begin{array}{l}
\mathrm{Ob}\,\uTy=\Ty=\ \text{free monoid on the two generators } \Delta^+,\ \Delta^- \text{ above.}\\[1ex]
\text{If } \tau,\tau'\in\Ty,\ \Hom(\tau,\tau')=\Hom_{\Cat}(I_\tau,I_{\tau'}).
\end{array}
\right.
\]

(\footnote[7]{A somewhat botched definition, one must restrict oneself to homomorphisms from $I_\tau$ to $I_{\tau'}$ compatible with source and target - cf. p.~29.}). We will note that we have not only a set $\Hom_{\Cat}(I_\tau,I_{\tau'})$, but even a category $\uHom(I_\tau,I_{\tau'})$, which here amounts to saying that $\Hom(I_\tau,I_{\tau'})$ (i.e. $\Hom(\tau,\tau')$) can be viewed as an \textit{ordered} set. But for the moment we are going to ignore this order relation.

\msp{page 26}
The composition of morphisms is the composition of functors between categories $I_\tau$. Thus we have a \textit{fully faithful} (\footnote[8]{false, cf. marginal annotation on the preceding page - we do not take \textit{all} the functors in Cat between $I_\tau$ and $I_{\tau'}$.}) canonical functor
\[
\tau \longmapsto I_\tau : \uTy \xrightarrow{\rho} (\Cat) \qquad \text{(\footnotemark[9]),}
\]
\footnotetext[9]{\textit{realization} functor of a (combinatorial) type of intervals.}% footnote text printed at the foot of p. 17
\printedpage{17}
% head: A. Grothendieck : Dérivateurs, Ch. VII, Édition des Universités de Montpellier II et Paris VII
% foot: Version 22 décembre 2001
by which the category $I_\tau$ \texttt{[rather} $\uTy$\texttt{]} could be identified with a full subcategory of (Cat) (that of the $I_\tau$). But I prefer not to identify $\tau\in\Ty$ (combinatorial datum) with $I_\tau$ (categorical datum).

Let $X$ be a (small?) category. We have a contrafunctor
\[
\left\{
\begin{array}{rcl}
S & \longmapsto & \uHom(S,X)\\[1ex]
(\Cat)^o & \longrightarrow & (\Cat)
\end{array}
\right.
\]
defined by $X$, whence by composing with $\rho$ a contrafunctor
\[
\uTy^{o} \xrightarrow{\mathrm{CH}_X} \Cat.
\]
This is the contrafunctor
\[
\left\{
\begin{array}{l}
\tau \longmapsto \uHom(I_\tau,X) \overset{\text{def}}{=} \uCh_\tau(X),\\[1ex]
\text{i.e. } \mathrm{CH}_X(\tau)=\uCh_\tau(X).
\end{array}
\right.
\]
But a contrafunctor on a small category $\uTy$ gives rise to a \textit{fibered category} (and even split) over it.

\msp{page 27}
I am going to denote this fibered category by $\uCh(X)$. (Provisionally - if it turns out to be too big, I will restrict it \dots) Thus we have
\[
\begin{tikzcd}[row sep=large]
\uCh(X) \arrow[d, "\pi_X"] \\
\uTy\ .
\end{tikzcd}
\]
The fiber of $\uCh(X)$ at $\tau\in\mathrm{Ob}\,\uTy$ is the category $\uCh_\tau(X)$ of paths of type $\tau$ in $X$. This therefore tells us what the arrows in $\uCh(X)$ over the identity arrows in $\uTy$ are. But what are the more general arrows? We give ourselves an arrow in $\uTy$
\[
\varphi : \tau' \longrightarrow \tau, \qquad \text{i.e. } I_{\tau'} \xrightarrow{\varphi} I_\tau,
\]
what are the arrows of $\uCh(X)$ over $\varphi$? By the general sorites of fibered categories, if
\[
\begin{array}{lclcl}
c\in\uCh(X)_\tau & = & \uCh_\tau(X) & \overset{\text{def}}{=} & \uHom(I_\tau,X)\\[1ex]
c'\in\uCh(X)_{\tau'} & = & \uCh_{\tau'}(X) & \overset{\text{def}}{=} & \uHom(I_{\tau'},X)
\end{array}
\]
we have
\[
\boxed{\Hom_\varphi(c',c)\simeq\Hom(c',\underbrace{\varphi^*(c)}_{=c\circ\varphi}).}
\]
Thus: the homomorphisms from the
\printedpage{18}
% head: A. Grothendieck : Dérivateurs, Ch. VII, Édition des Universités de Montpellier II et Paris VII
% foot: Version 22 décembre 2001
\msp{page 28}
path $c'$ of type $\tau'$ in $X$, to the path $c$ of type $\tau$, ``over'' the morphism $\varphi$ from $\tau'$ to $\tau$, i.e. the functor
\[
\varphi : I_{\tau'} \longrightarrow I_\tau,
\]
is a homomorphism $v$ from $c'$ to $c\circ\varphi$, i.e.
\[
\begin{tikzcd}[row sep=large]
I_{\tau'} \arrow[r, "\varphi"] \arrow[dr, "c'"', ""{name=A}] & I_\tau \arrow[d, "c", ""'{name=B}] \\
& X.
\arrow[from=A, to=B, bend left=20, shorten <=6pt, shorten >=2pt, "\circled{v}"]
\end{tikzcd}
\]
The composition of arrows is obvious:
\[
\begin{tikzcd}[row sep=large]
I_{\tau''} \arrow[r, "\varphi'"] \arrow[drr, "c''"', ""{name=C}] & I_{\tau'} \arrow[r, "\varphi"] \arrow[dr, "c'"', ""{name=A}] & I_\tau \arrow[d, "c", ""'{name=B}] \\
& & X,
\arrow[from=A, to=B, bend left=20, shorten <=6pt, shorten >=2pt, "v"]
\arrow[from=C, to=A, shorten <=4pt, shorten >=4pt, "v'"]
\end{tikzcd}
\]
we have
\[
\begin{array}{rccl}
v : & c' & \longrightarrow & c\circ\varphi, \quad \text{whence } c'\circ\varphi' \xrightarrow{v*\varphi'} \underbrace{(c\circ\varphi)\circ\varphi'}_{=c\circ(\varphi\varphi')}\\[3ex]
v' : & c'' & \longrightarrow & c'\circ\varphi'
\end{array}
\]
and
\[
v''=(v*\varphi')\circ v', \qquad \text{thus} \qquad
\begin{tikzcd}
c'' \arrow[r, "v'"] \arrow[rr, bend right=12, "\circled{v''}"'] & c'\circ\varphi' \arrow[r, "v*\varphi'"] & c\circ(\varphi\varphi').
\end{tikzcd}
\]
But do we have source and target functors
\[
\uCh(X) \overset{\sigma}{\underset{\beta}{\rightrightarrows}} X \qquad
\begin{array}{ccl}
c & \longmapsto & \text{source}(c)\\[1ex]
c & \longmapsto & \text{target}(c)
\end{array}
\qquad ?
\]
It would therefore be necessary for a homomorphism $v$ as above
\[
v : c' \longrightarrow c\circ\varphi
\]
to define
\[
\begin{array}{lrcl}
\sigma(c) : & \text{source}(c') & \longrightarrow & \text{source}(c)\\[1ex]
\beta(c) : & \text{target}(c') & \longrightarrow & \text{target}(c).
\end{array}
\]